Modern LLMs usually employ a two-stage inference process: \textbf{Prefill Stage} and \textbf{Decode Stage}. The Prefill Stage serves as the initial step where the key-value cache (KV cache) for the prompt (also called prefix) sequences are built for each KV head in the Transformer architecture\footnote{Note that here we discriminate KV head from attention head because of Grouped Query Attention (GQA)}. Given hidden states $\mathbf{X}\in\mathbb{R}^{l\times d_{\mathrm{model}}}$, each attention layer computes $\mathbf{Q}=\mathbf{X}\mathbf{W}_q\in\mathbb{R}^{l\times h_q\times d_h}$, $\mathbf{K}=\mathbf{X}\mathbf{W}_k\in\mathbb{R}^{l\times h_{\mathrm{kv}}\times d_h}$, and $\mathbf{V}=\mathbf{X}\mathbf{W}_v\in\mathbb{R}^{l\times h_{\mathrm{kv}}\times d_h}$ after reshaping, where $h_q$ is the number of query heads, $h_{\mathrm{kv}}$ is the number of KV heads, and $d_h$ is the head dimension. After prefilling, all the key-value pairs are preserved in the memory so that they can be reused during the decoding of each token. This will result in two vectors $\textbf{K}_{\text{cache}}=[\textbf{K}_1;\ldots;\textbf{K}_l]\in \mathbb{R}^{l\times h_{\mathrm{kv}}\times d_h}$ and $\textbf{V}_{\text{cache}}=[\textbf{V}_1;\ldots;\textbf{V}_l]\in \mathbb{R}^{l\times h_{\mathrm{kv}}\times d_h}$. The GPU High Bandwidth Memory (HBM) typically needs to spare a significant amount of spaces to store these KV cache. For instance, for \texttt{Qwen/Qwen3-32B} in FP16 precision, it will take up $16$ GB to save a sequence of $64$K tokens. In contrast, \texttt{Qwen/Qwen3-4B} requires roughly half of the space to save these amount of KV cache\footnote{\url{https://lmcache.ai/kv_cache_calculator.html} provides calculations for some other models}. The Prefill Stage is usually \emph{compute-bound}, since it involves matrix-matrix multiplications that saturates the GPU's arithmetic units. We use \emph{Time to First Token} (TTFT) to calculate the time spent from waiting the first token to emit to represent the prefilling latency.

Once the prompt is processed, the model enters the Decode Stage. For each step $t$, the model processes the current prefix and concatenates a new KV pair: $\textbf{K}_{\text{cache}}^{(t)}=[\textbf{K}_{\text{cache}}^{(t-1)};\textbf{K}^{(t)}],\textbf{V}_{\text{cache}}^{(t)}=[\textbf{V}_{\text{cache}}^{(t-1)};\textbf{V}^{(t)}]$. 
% \begin{align*}
%     &\textbf{K}_{\text{cache}}^{(t)}=[\textbf{K}_{\text{cache}}^{(t-1)};\textbf{K}^{(t)}] \\
%     &\textbf{V}_{\text{cache}}^{(t)}=[\textbf{V}_{\text{cache}}^{(t-1)};\textbf{V}^{(t)}]
% \end{align*}
And at the same time emits the output token. Different from Prefill Stage, Decode Stage is memory-bound since at each decoding step, the GPU spends more time waiting for the KV cache to arrive from HBM. Consequently, the decoding speed is usually bottlenecked by the memory bandwidth.
When the HBM is saturated, part of the KV cache will be either emitted or offloaded to CPU~\cite{jin2024compute}.
% Since CPU bandwidth is linear and does not incur quadratic compute in long-context prefilling, \emph{CPU offload} is often a better strategy which results in significantly cheaper prefilling costs.
CPU offload is often a better strategy which results in linear prefilling costs compared to the quadratic compute in long-context prefilling.
Typically, we measure decoding efficiency through \emph{Time per Output Token} (TPOT), that tracks the average generation speed per token, excluding the initial prefilling costs.